
\documentclass[11pt,a4paper]{article}
\usepackage[a4paper,margin=2.05cm]{geometry}
\usepackage[spanish,es-noquoting]{babel}
\usepackage[T1]{fontenc}
\usepackage[utf8]{inputenc}
\usepackage{lmodern,microtype}
\usepackage{amsmath,amssymb,amsthm,mathtools}
\usepackage{xcolor,hyperref}
\pagecolor{white}\color{black}
\setlength{\parindent}{0pt}
\setlength{\parskip}{0.65em}
\title{\bfseries N62 --- Derivación de la firma de frontera \(R_{36}\) desde \(B_4\)}
\author{HMT / auditoría técnica}
\date{}
\begin{document}
\maketitle
\section*{Resultado}
La frontera anterior es
\[
u_4^\pi=110221,\qquad u_4^e=102011,\qquad u_4^\varphi=010200.
\]
Se verifica:
\[
\boxed{u_5^\varphi=u_4^e=102011.}
\]
Y:
\[
\boxed{u_5^\pi+u_5^e=u_4^\varphi A_WL_1^3=210112.}
\]
Por tanto:
\[
\boxed{q=012120,\qquad a=111101.}
\]
Bajo saturación de carry, la ocupación se lee como
\[
\operatorname{colw}_j=2+\mathbf1_{\{u^\varphi_j\ne0,\ (u^\pi+u^e)_j\ne2\}},
\]
lo que da
\[
\boxed{\operatorname{colw}=223232.}
\]
La quinta triada decimal deja \(38\cdot150\cdot150=855000\) ternas. Con \(r=102\), la firma y la poda seleccionan de forma única:
\[
\boxed{(u_5^\pi,u_5^e,u_5^\varphi)=(222220,021222,102011).}
\]
\section*{Estado}
\[
\boxed{q,a,\operatorname{colw}\text{ quedan derivados.}}
\]
\[
\boxed{r=102\text{ no queda forzado por }q,a,c,\operatorname{colw};\text{ es la carga de fila residual.}}
\]
\end{document}
